\documentclass[UTF8,a4paper,11pt]{ctexart}

% ── Page geometry ──
\usepackage[margin=2.2cm,top=2.5cm,bottom=2.5cm]{geometry}

% ── Typography ──
\usepackage{fontspec}
\usepackage{setspace}
\setstretch{1.2}

% ── Graphics ──
\usepackage{graphicx}
\graphicspath{{/home/zhiwen/projects/mooc2handout-skill/coursera-mcp/module-1/frames/}}

% ── Colors ──
\usepackage{xcolor}
\definecolor{accent}{HTML}{1a73e8}
\definecolor{darkgray}{HTML}{333333}
\definecolor{lightgray}{HTML}{f5f5f5}
\definecolor{keyframebg}{HTML}{fafafa}
\definecolor{videoblue}{HTML}{065fd4}
\definecolor{notebg}{HTML}{e8f0fe}
\definecolor{warnbg}{HTML}{fef7e0}

% ── Links ──
\usepackage{hyperref}
\hypersetup{
  colorlinks=true,
  linkcolor=accent,
  urlcolor=videoblue,
  citecolor=accent,
  pdftitle={模型上下文协议导论},
  pdfauthor={}
}

% ── Layout ──
\usepackage{enumitem}
\setlist[itemize]{leftmargin=1.8em,itemsep=2pt}
\setlist[enumerate]{leftmargin=1.8em,itemsep=2pt}
\usepackage{booktabs}
\usepackage{longtable}
\usepackage{tcolorbox}
\usepackage{titlesec}
\usepackage{fancyhdr}
\usepackage{lastpage}
\usepackage{amssymb}

% ── Header/Footer ──
\pagestyle{fancy}
\fancyhf{}
\fancyhead[L]{\small\textcolor{darkgray}{模型上下文协议导论}}
\fancyhead[R]{\small\textcolor{darkgray}{第一单元：认识 MCP}}
\fancyfoot[C]{\small\textcolor{darkgray}{\thepage\ / \pageref{LastPage}}}
\renewcommand{\headrulewidth}{0.4pt}

% ── Section styling ──
\titleformat{\section}
  {\Large\bfseries\color{accent}}
  {\thesection}{1em}{}
\titleformat{\subsection}
  {\large\bfseries\color{darkgray}}
  {\thesubsection}{1em}{}

% ── Custom environments ──
\newtcolorbox{keyconcept}{
  colback=notebg, colframe=accent,
  fonttitle=\bfseries, boxrule=0.5pt,
  left=8pt, right=8pt, top=6pt, bottom=6pt
}
\newtcolorbox{supplement}{
  colback=lightgray, colframe=darkgray,
  fonttitle=\bfseries, boxrule=0.3pt,
  left=8pt, right=8pt, top=6pt, bottom=6pt
}
\newtcolorbox{exercise}{
  colback=warnbg, colframe=orange,
  fonttitle=\bfseries, boxrule=0.5pt,
  left=8pt, right=8pt, top=6pt, bottom=6pt
}

% ── Video link command ──
\newcommand{\videolink}[2]{%
  \href{#1}{\textcolor{videoblue}{\textbf{#2}}}%
}

% ── Keyframe figure command ──
\newcommand{\keyframe}[2]{%
  \begin{center}
    \fcolorbox{lightgray}{keyframebg}{%
      \includegraphics[width=0.75\linewidth]{#1}%
    }
    \par\vspace{2pt}
    {\small\textcolor{darkgray}{#2}}
  \end{center}
}

\begin{document}

\begin{center}
  \vspace*{2cm}
  {\Huge\bfseries\color{accent} 模型上下文协议导论}\\[0.3cm]
  {\small\textcolor{darkgray}{Introduction to Model Context Protocol}}\\[0.8cm]
  {\LARGE 第一单元：认识 MCP}\\[0.3cm]
  {\small\textcolor{darkgray}{Module 1: Getting Started with MCP}}\\[1.2cm]
  {\large 授课教师： }\\[0.5cm]
  {\large\textcolor{darkgray}{由 mooc2handout 自动生成}}\\[0.3cm]
  {\small\textcolor{darkgray}{\today}}
  \vfill
  {\small 本讲义由课程字幕、补充材料与可选的 AI 推断关键帧自动生成。}
\end{center}
\newpage

\tableofcontents
\newpage

\section*{本单元知识地图}
\addcontentsline{toc}{section}{本单元知识地图}
\begin{longtable}{p{3.1cm}p{10.2cm}}
\toprule
主题 & 核心问题 \\
\midrule
MCP 是什么 & 它如何把模型和外部能力连接起来？ \\
Client / Server & 谁负责通信，谁负责暴露能力？ \\
三类原语 & tools、resources、prompts 各自解决什么问题？ \\
\bottomrule
\end{longtable}

\section*{0. 概要}
\addcontentsline{toc}{section}{0. 概要}
这一单元先回答一个最根本的问题：为什么需要 MCP。
它把“模型如何安全、稳定地接入外部能力”变成一个标准协议，
也把应用、client、server 的职责分开，避免把 API 调用直接硬塞进提示词里。

\begin{itemize}
  \item 理解 MCP 解决的是“连接层”问题，而不是替代模型本身。
  \item 认识 client / server 的分工。
  \item 初步建立 tools、resources、prompts 这三类原语的直觉。
\end{itemize}

\section*{1. 欢迎与课程路线图}
\addcontentsline{toc}{section}{1. 欢迎与课程路线图}
\noindent\textbf{回看：}\videolink{https://www.coursera.org/learn/introduction-to-model-context-protocol/home/module/1}{欢迎与课程路线图}

\textbf{本讲重点：} 先建立课程地图，再明确学习前提。

这一讲不是先写代码，而是先把整门课的路线图铺开：你会先理解 MCP 要解决什么问题，
再拆开 client 和 server 的职责，最后进入 tools、resources、prompts 三类服务器能力。

讲者同时交代了学习前提：需要基础 Python、本地运行环境，以及后面会用到的 uv-cli。
把环境准备好，后面的动手部分才不会被琐碎问题打断。

\begin{itemize}
  \item 先知道这门课要解决什么，再进入实现细节。
  \item 本课程会把概念、协议和动手项目结合起来。
  \item 准备好 Python 和 uv-cli 是进入实操阶段的前提。
\end{itemize}

\begin{keyconcept}
\textbf{自动摘要}
根据字幕自动扩写，这一讲再补充几条最值得记住的线索。
\begin{itemize}
  \item \textbf{形式定义}: 这一讲梳理 AND / OR / NOT / IF...THEN 等连接词的真值条件，并提醒不要把日常语言的直觉直接搬进数学里。
\end{itemize}
\end{keyconcept}

\section*{2. 认识 MCP}
\addcontentsline{toc}{section}{2. 认识 MCP}
\noindent\textbf{回看：}\videolink{https://www.coursera.org/learn/introduction-to-model-context-protocol/home/module/1}{认识 MCP}

\textbf{本讲重点：} MCP 是连接模型与外部能力的标准通信层。

这一讲用 GitHub 聊天机器人的例子说明：如果每次都手写 API 调用、拼接上下文、处理权限，
工程会迅速变得凌乱。MCP 的价值就是把这些重复工作抽象成统一协议。

在这个框架里，server 不只是“一个后端”，它还可以把 tools、resources、prompts
统一暴露出来，让应用和模型通过标准方式协作。

\begin{itemize}
  \item tools 负责让模型执行明确动作，比如查询 Pull Request。
  \item resources 负责把数据送进上下文。
  \item prompts 负责把可复用的指令模板交给应用或用户。
\end{itemize}

\begin{keyconcept}
\textbf{自动摘要}
根据字幕自动扩写，这一讲再补充几条最值得记住的线索。
\begin{itemize}
  \item \textbf{形式定义}: 这一讲梳理 AND / OR / NOT / IF...THEN 等连接词的真值条件，并提醒不要把日常语言的直觉直接搬进数学里。
\end{itemize}
\end{keyconcept}

\section*{3. MCP 客户端}
\addcontentsline{toc}{section}{3. MCP 客户端}
\noindent\textbf{回看：}\videolink{https://www.coursera.org/learn/introduction-to-model-context-protocol/home/module/1}{MCP 客户端}

\textbf{本讲重点：} client 是通信入口，transport 可以有多种实现。

MCP 并不强制一种传输协议：同机时可以走 stdin/stdout，分布式时也可以走 HTTP 或
WebSockets。客户端负责把应用请求翻译成协议消息，再把 server 的响应交回上层。

这一讲还点出了几个关键消息类型，比如 \texttt{ListToolsRequest}、
\texttt{ListToolsResult} 和 \texttt{CallToolRequest}。理解这些消息，
才能看懂 client 和 server 之间到底交换了什么。

\begin{itemize}
  \item client 不是能力本身，而是访问 server 的入口。
  \item 协议定义了消息类型和交互方式。
  \item 同一个 server 可以被不同应用复用。
\end{itemize}

\begin{keyconcept}
\textbf{自动摘要}
根据字幕自动扩写，这一讲再补充几条最值得记住的线索。
\begin{itemize}
  \item \textbf{形式定义}: 这一讲梳理 AND / OR / NOT / IF...THEN 等连接词的真值条件，并提醒不要把日常语言的直觉直接搬进数学里。
\end{itemize}
\end{keyconcept}

\section*{4. 本单元最该记住的五句话}
\addcontentsline{toc}{section}{4. 本单元最该记住的五句话}
\begin{enumerate}
  \item MCP 是一个标准通信层，不是单一模型或框架。
  \item Client 负责连接与转发，Server 负责暴露能力。
  \item Tools、resources、prompts 是 server 的三类原语。
  \item 传输层可以是 stdio、HTTP、WebSockets 等。
  \item 先理解协议，再写代码，会更容易看懂后面的项目。
\end{enumerate}

\section*{5. 练习题}
\addcontentsline{toc}{section}{5. 练习题}
\subsection*{A. 选择题（5题）}
\begin{enumerate}
  \item MCP 最核心的作用是什么？\\
  A. 代替 Python 解释器 \quad B. 让模型和外部能力用统一协议通信 \quad C. 生成 UI 组件 \quad D. 自动训练模型

  \item MCP server 通常暴露哪三类能力？\\
  A. 变量、类、函数 \quad B. 文件、网络、数据库 \quad C. tools、resources、prompts \quad D. 终端、浏览器、IDE

  \item 课程示例中的应用主要想让模型访问什么？\\
  A. GitHub 数据 \quad B. 本地相册 \quad C. 股票行情 \quad D. 邮件附件

  \item 为什么说 MCP 是 transport agnostic？\\
  A. 只能用 HTTP \quad B. 只支持同机进程 \quad C. 可通过多种协议通信 \quad D. 必须走数据库

  \item ListToolsRequest 的作用是什么？\\
  A. 创建工具 \quad B. 请求服务器列出可用工具 \quad C. 发送提示词 \quad D. 下载资源

\end{enumerate}

\subsection*{B. 判断题（3题）}
\begin{enumerate}
  \item Client 负责把服务器能力暴露给模型。（\quad）
  \item 同一机器上的 client 和 server 可以通过 stdin/stdout 通信。（\quad）
  \item MCP 不需要明确的消息类型。（\quad）
\end{enumerate}

\subsection*{C. 简答题（2题）}
\begin{enumerate}
  \item 用自己的话解释 client 和 server 的分工。
  \item 为什么 GitHub 聊天机器人适合用 MCP 来实现？
\end{enumerate}

\section*{6. 参考答案}
\addcontentsline{toc}{section}{6. 参考答案}
\subsection*{选择题答案}
1.B \quad 2.C \quad 3.A \quad 4.C \quad 5.B

\subsection*{判断题答案}
1. 错 \quad 2. 对 \quad 3. 错

\subsection*{简答题要点}
\begin{enumerate}
  \item Client 负责建立连接、发送请求和接收响应；Server 负责把工具、资源和提示词等能力暴露出来。
  \item 因为 GitHub 场景需要稳定地访问外部 API 和上下文，MCP 可以把这些能力标准化并复用。
\end{enumerate}

\end{document}